%! AP Statistics — Unit 1, Lesson 1.8 — Group Activity
\documentclass[10pt]{article}
\usepackage{apstats-article}
\usepackage{apstats-boxes}

%=============================================================================
\begin{document}

\pageheader{Unit 1, Lesson 1.8}{Group Activity: Boxplots \& Modified Boxplots}
\namepartnerperiod

\vspace{0.10in}

\begin{tcolorbox}[colback=sky, colframe=navy, boxrule=0.8pt, arc=2mm,
    left=3mm, right=3mm, top=2mm, bottom=2mm]
\small
\textbf{Part 1} (group, 8 min): Compute the five-number summary and check for
outliers.\\
\textbf{Part 2} (group, 7 min): Interpret a pre-drawn modified boxplot.\\
\textbf{Part 3} (individual, 5 min): Write a SOCS description and compare
displays.
\end{tcolorbox}

\vspace{0.10in}

% ── PART 1 ───────────────────────────────────────────────────────────────────
{\large\bfseries\color{navy} Part 1 \quad Compute the Five-Number Summary}

\vspace{0.10in}

\noindent\small A school nurse measured the resting pulse rates (beats per minute)
of 17 student athletes. The sorted data are:

\vspace{0.08in}
\begin{center}
\small
52,\enspace 56,\enspace 58,\enspace 60,\enspace 62,\enspace 63,\enspace 64,\enspace
66,\enspace 68,\enspace 70,\enspace 72,\enspace 73,\enspace 75,\enspace 78,\enspace
80,\enspace 84,\enspace 110
\end{center}

\vspace{0.12in}

\begin{enumerate}[leftmargin=1.5em, itemsep=8pt]

  \item How many data values are there? Which position is the median?\\[5pt]
        $n = \blank{1.5cm}$ values. \quad The median is the \blank{2cm} value.
        \quad $M = \blank{2cm}$ bpm

  \item Find $Q_1$ (median of the lower half) and $Q_3$ (median of the upper half).\\[5pt]
        Lower half (8 values): \blank{8cm}\\[4pt]
        $Q_1 = \blank{2.5cm}$ bpm \hfill
        Upper half (8 values): \blank{8cm}\\[4pt]
        $Q_3 = \blank{2.5cm}$ bpm \hfill
        $\text{IQR} = Q_3 - Q_1 = \blank{3cm}$ bpm

  \item Record the five-number summary:\\[5pt]
        \renewcommand{\arraystretch}{1.5}
        \begin{tabularx}{0.85\linewidth}{@{} Y Y Y Y Y @{}}
        \rowcolor{sky}
        $\min$ & $Q_1$ & $M$ & $Q_3$ & $\max$ \\
        \midrule
        \blank{1.5cm} & \blank{1.5cm} & \blank{1.5cm} & \blank{1.5cm} & \blank{1.5cm} \\
        \end{tabularx}

  \item Calculate the lower and upper fences. Is 110 an outlier? Show work.\\[5pt]
        Lower fence $= Q_1 - 1.5(\text{IQR}) = \blank{7cm}$\\[4pt]
        Upper fence $= Q_3 + 1.5(\text{IQR}) = \blank{7cm}$\\[4pt]
        Is 110 an outlier? \blank{2cm} \quad Because: \blank{6cm}

  \item What is the last \textit{non-outlier} value on the right? \blank{3cm} bpm.\\
        This is where the right whisker of a modified boxplot would end.

\end{enumerate}

\vspace{0.12in}

% ── PART 2 ───────────────────────────────────────────────────────────────────
{\large\bfseries\color{navy} Part 2 \quad Interpret the Modified Boxplot}

\vspace{0.10in}

\noindent\small The modified boxplot below was made from the pulse rate data.

\vspace{0.10in}

\begin{center}
\begin{tikzpicture}[scale=1.1]
  % Axis: 45 to 120, scale (v-45)*0.10; label every 10
  % Ticks at 50,60,70,80,90,100,110,120
  \draw[thick] (0.4,0) -- (8.0,0);
  \foreach \v/\lbl in {0.50/50, 1.50/60, 2.50/70, 3.50/80,
                        4.50/90, 5.50/100, 6.50/110, 7.50/120} {
    \draw (\v,-0.12) -- (\v,0.12);
    \node[below, font=\small] at (\v,-0.12) {\lbl};
  }
  \node[below, font=\small\itshape] at (4.0,-0.54) {Pulse rate (bpm)};

  % scale: (v-45)*0.10
  % min=52: x=0.70, Q1=62.5: x=1.75, M=68: x=2.30,
  % Q3=76.5: x=3.15, last non-outlier=84: x=3.90, outlier=110: x=6.50
  \draw[thick, navy] (0.70, 0.35) -- (1.75, 0.35);   % left whisker min→Q1
  \draw[thick, navy] (1.75, 0.10) rectangle (3.15, 0.60);  % box Q1→Q3
  \draw[thick, navy] (2.30, 0.10) -- (2.30, 0.60);   % median
  \draw[thick, navy] (3.15, 0.35) -- (3.90, 0.35);   % right whisker Q3→last non-outlier
  \draw[thick, navy] (0.70, 0.20) -- (0.70, 0.50);   % left cap
  \draw[thick, navy] (3.90, 0.20) -- (3.90, 0.50);   % right cap
  \filldraw[navy] (6.50, 0.35) circle (0.10);         % outlier dot at 110
\end{tikzpicture}
\end{center}

\vspace{0.10in}

\begin{enumerate}[leftmargin=1.5em, itemsep=8pt, resume]

  \item From the boxplot, identify: $Q_1 = \blank{2cm}$, $M = \blank{2cm}$,
        $Q_3 = \blank{2cm}$. Do these match your calculations?\\[4pt]
        \writeline

  \item What is the value of the outlier dot? \blank{2cm} bpm. \quad
        What does the right whisker end represent? \blank{5cm}

  \item Approximately what percent of the athletes had pulse rates between
        62.5 and 76.5 bpm? \blank{3cm}\%. \quad
        Between 52 and 68 bpm? \blank{3cm}\%.

\end{enumerate}

\vspace{0.12in}

% ── PART 3 ───────────────────────────────────────────────────────────────────
{\large\bfseries\color{navy} Part 3 \quad SOCS Description \& Comparing Displays}

\vspace{0.10in}

\begin{enumerate}[leftmargin=1.5em, itemsep=8pt, resume]

  \item Write a full \textbf{SOCS} description of the pulse rate distribution
        based on the modified boxplot. Use the context (``pulse rates of student
        athletes'') in every sentence.\\[5pt]
        \textbf{Shape:} \writeline\\[1pt]\writeline\\[5pt]
        \textbf{Outlier(s):} \writeline\\[1pt]\writeline\\[5pt]
        \textbf{Center:} \writeline\\[5pt]
        \textbf{Spread:} \writeline

  \item Name \textbf{one thing} the modified boxplot shows clearly that a
        histogram of the same data would \textit{not} show as clearly.\\[5pt]
        \writeline

  \item Name \textbf{one thing} a histogram would reveal that the boxplot
        cannot.\\[5pt]
        \writeline

\end{enumerate}

\end{document}
